\documentclass[conference]{IEEEtran}
\IEEEoverridecommandlockouts
\usepackage{cite}
\usepackage{amsmath,amssymb,amsfonts}
\usepackage{algorithmic}
\usepackage{graphicx}
\usepackage{textcomp}
\usepackage{xcolor}
\usepackage{hyperref}
\usepackage{listings}
\usepackage{tabularx}
\usepackage{balance}

\newcolumntype{C}{>{\centering\arraybackslash}X}

\usepackage{tikz}
\usetikzlibrary{shapes,arrows,positioning,fit,backgrounds}

\def\BibTeX{{\rm B\kern-.05em{\sc i\kern-.025em b}\kern-.08em
    T\kern-.1667em\lower.7ex\hbox{E}\kern-.125emX}}

\lstset{
  basicstyle=\ttfamily\scriptsize,
  columns=fullflexible,
  frame=single,
  breaklines=true,
  captionpos=b,
}

\begin{document}

\title{Zero-Trust Architecture for MCP-Based AI Agents: Continuous Identity and Behavioral Monitoring}

\author{\IEEEauthorblockN{Yoshihiro Ando}
\IEEEauthorblockA{\textit{Independent Researcher} \\
yosh5295@gmail.com \\
Tokyo, Japan}
}

\maketitle

\begin{abstract}
As autonomous AI agents evolve into distributed system operators via the Model Context Protocol (MCP), the assumption of network-level security becomes obsolete. This paper, representing \textbf{Phase 2: Zero Trust (Execution Security)} of a comprehensive security trilogy, presents a framework for securing agentic workloads through "Spatial Security" and continuous behavioral monitoring. Building upon the structural guardrails of Phase 1, we implement a Zero Trust Architecture (ZTA) centered on asymmetric workload identity (RSA) and a \textbf{Mamba-based Reasoning Sentinel}. The sentinel utilizes a NumPy-only State Space Model (SSM) to perform real-time anomaly detection on the agent's reasoning traces, providing a behavioral Intrusion Detection System (IDS) for AI logic. Our results demonstrate that the combination of sub-millisecond cryptographic verification and lightweight behavioral monitoring effectively mitigates risks like prompt injection and semantic drift without impeding agentic responsiveness.
\end{abstract}

\begin{IEEEkeywords}
Zero Trust Architecture, AI Agents, Model Context Protocol, Workload Identity, Mamba Sentinel, Intrusion Detection System (IDS), Reasoning Integrity.
\end{IEEEkeywords}

\section{Introduction}
Phase 1 of this trilogy \cite{ando2026guardrails} established structural containment via sandboxing and "No-Shell" architectures. However, a "jailbroken" agent might still perform harmful actions that are syntactically valid but semantically malicious. For instance, an agent permitted to "manage documents" might be tricked into deleting a critical database file if it can be represented as a document. This necessitates a shift from static containment to dynamic, behavioral verification.

We propose a Zero Trust Architecture (ZTA) for AI agents, where the principle of "Never Trust, Always Verify" is applied not just to the network, but to the \textit{intent} and \textit{identity} of the AI's reasoning process.

This paper is the second installment of our security trilogy:
\begin{itemize}
    \item \textbf{Phase 1: Guardrails (Pre-execution Security) \cite{ando2026guardrails}.} Structural containment and sandboxing.
    \item \textbf{Phase 2: Zero Trust (Execution Security).} The current work, focusing on \textbf{Continuous Verification} via identity propagation and behavioral integrity.
    \item \textbf{Phase 3: Post-Quantum (Temporal Security) \cite{ando2026pqc}.} Long-term integrity and quantum-resistant governance.
\end{itemize}

\section{Asymmetric Identity and Tool Orchestration}
In a Zero Trust model, an agent's authority must be tied to a verifiable identity rather than its session or location.

\subsection{Workload Identity (RSA-2048)}
\texttt{llm-cli} assigns each agent instance a unique RSA-2048 key pair. When the agent requests a tool call via MCP, the system generates a signed \textbf{Identity Token (JWT)}. This token acts as a "Passport" for the agent, containing:
\begin{itemize}
    \item \textbf{sub (Subject):} The agent's unique ID (e.g., \texttt{user@hostname}).
    \item \textbf{roles:} Assigned RBAC roles (e.g., \texttt{admin}, \texttt{user}).
    \item \textbf{iat/exp:} Issued-at and expiration timestamps to prevent replay attacks.
\end{itemize}
This ensures that remote MCP servers can cryptographically verify that the tool call originated from a legitimate, untampered agent instance.

\subsection{Role-Based Access Control (RBAC)}
Permissions are decoupled from the LLM's prompt. Even if the LLM "believes" it has root access due to a prompt injection, the underlying security engine enforces a strict RBAC policy. If a tool (e.g., \texttt{edit\_file}) requires the \texttt{admin} role, the execution is blocked regardless of the AI's "monologue" unless the signed identity token carries that claim.

\section{Behavioral IDS: The Mamba Sentinel}
While identity ensures *who* is calling the tool, it does not ensure *why*. To detect semantic attacks like prompt injection, we introduce the \textbf{Mamba Sentinel}.

\subsection{The Need for Lightweight Monitoring}
Traditional monitoring using secondary LLMs (e.g., using GPT-4o-mini to verify GPT-4o's output) is slow (1-2 seconds of latency) and expensive. We utilize a \textbf{Mamba (State Space Model)} implementation \cite{mamba2023} optimized with pure NumPy. Unlike Transformers, Mamba's linear scaling and recurrent state make it ideal for monitoring streaming reasoning traces at a byte level with sub-millisecond overhead.

\subsection{Reasoning Integrity and Surprise Score}
The Sentinel monitors the LLM's "internal monologue" (reasoning tokens) in real-time. It calculates a \textbf{Surprise Score} based on the cross-entropy loss between the Sentinel's predicted next-token distribution and the agent's actual output.

\begin{equation}
L_{sentinel} = -\sum P_{sentinel}(x_t | h_{t-1}) \log P_{agent}(x_t)
\end{equation}

A sudden spike in "surprise" indicates a statistical anomaly or "Semantic Drift." This often happens when the agent transitions from a helpful assistant state to an adversarial state (e.g., suddenly generating obfuscated code or trying to access sensitive files after a prompt injection).

\subsection{Dynamic Intervention (Green/Yellow/Red)}
The system maintains a real-time status:
\begin{itemize}
    \item \textbf{Green:} Normal operation.
    \item \textbf{Yellow:} Minor anomaly detected. The system increases the entropy-scanning sensitivity and alerts the user.
    \item \textbf{Red:} Significant anomaly. The system \textbf{automatically suspends execution} and forces a Human-in-the-Loop (HITL) review.
\end{itemize}

\begin{figure}[htbp]
\centering
\resizebox{\columnwidth}{!}{
    \input{trilogy\_phase2_architecture.tex}
}
\caption{Phase 2 Architecture: Continuous verification via Asymmetric Identity and Mamba Behavioral Sentinel.}
\label{fig:zt_arch}
\end{figure}

\section{Evaluation}
We evaluated the Zero Trust framework on an \textbf{AMD Ryzen 5 8540U} system.

\subsection{Execution Latency (Empirical)}
Table \ref{tab:zt_perf} summarizes the latencies. The Mamba Sentinel adds only ~10ms to the reasoning process, which is nearly two orders of magnitude faster than a "Dual-LLM" verifier approach.

\begin{table}[htbp]
\caption{Zero Trust and IDS Execution Latency}
\begin{center}
\begin{tabular}{|l|c|c|}
\hline
\textbf{Component} & \textbf{Latency (ms)} & \textbf{Type} \\
\hline
Identity Verification (RS256) & 0.0296 & Cryptographic \\
Mamba Sentinel (SSM) & 10.3122 & ML-Inference \\
\hline
\end{tabular}
\label{tab:zt_perf}
\end{center}
\end{table}

\subsection{Detection Accuracy}
In testing with known prompt injection techniques (e.g., "DAN", "Ignore previous instructions"), the Mamba Sentinel detected escapes by identifying statistical anomalies in the reasoning trace before a harmful tool call was even attempted. This provides a "Defense in Depth" that complements the structural guardrails of Phase 1.

\section{Conclusion}
By decoupling reasoning from authority through asymmetric identities and providing real-time behavioral monitoring via the Mamba Sentinel, we ensure that autonomous agents remain "Verified by Identity and Monitored by Logic." This fulfills the promise of agentic productivity with robust, low-latency execution security. The final phase of our trilogy addresses resilience against future quantum threats and remote integrity proofs.

\begin{thebibliography}{00}
\bibitem{mamba2023} A. Gu and T. Dao, ``Mamba: Linear-Time Sequence Modeling with Selective State Spaces,'' 2023.
\bibitem{ando2026guardrails} Y. Ando, ``Secure-by-Default Guardrails for MCP-Based Tool Use in Multi-Modal LLM Agents,'' \textit{TechRxiv}, 2026.
\bibitem{ando2026pqc} Y. Ando, ``Post-Quantum Workload Identity and Long-Term Audit Integrity for MCP-Based AI Agents,'' \textit{TechRxiv}, 2026.
\end{thebibliography}

\end{document}
